\documentclass[a4paper,10pt,brazil]{article}
\usepackage{provastyle}
\usepackage{menukeys}

\begin{document}
\makeexamtitlepage{UFOP}{BCC222: Programação Funcional}{José Romildo}{2026/1}{Primeira Prova}{06}
\begin{multicols}{2}
\questionTitle{1}{Ação recursiva sem caso base não produz lista parcial}
\begin{flushright}
  \footnotesize
  \textit{\textbf{10: Ações de E/S recursivas}}\\
  \textit{c10-bordas-robustez-02}\\

\end{flushright}
Considere a ação:

\begin{haskellcode}
leParaSempre :: IO [Integer]
leParaSempre = do
  x <- readLn
  xs <- leParaSempre
  return (x:xs)
\end{haskellcode}

Qual alternativa descreve melhor o problema dessa definição como ação de E/S?

\begin{enumerate}
  \item Ela falha por erro de tipo, pois \haskellinline|x:xs| nunca pode ser usado dentro de \haskellinline|IO|.

  \item Ela é segura porque listas em Haskell são preguiçosas, então o programa pode usar o primeiro elemento antes de ler os demais.

  \item Ela imprime cada número lido automaticamente, mesmo sem chamada a \haskellinline|print|.

  \item Ela sempre produz a lista vazia, pois \haskellinline|return| ignora todos os valores lidos.

  \item Como não há caso base, a ação precisa continuar executando leituras indefinidamente antes de produzir o valor \haskellinline|IO [Integer]| esperado.


\end{enumerate}

\questionTitle{2}{Uso seguro de fromJust}
\begin{flushright}
  \footnotesize
  \textit{\textbf{6: Estruturas de dados básicas}}\\
  \textit{c06-maybe-operacoes-02}\\

\end{flushright}
Em qual situação o uso de \haskellinline{fromJust} é mais justificável neste capítulo?

\begin{enumerate}
  \item Apenas quando o conteúdo for uma string, mas não quando for um número.

  \item Sempre que a expressão aparecer dentro de uma lista.

  \item Sempre que quisermos evitar o uso de \haskellinline{fromMaybe}.

  \item Sempre que a expressão tiver o tipo \haskellinline{Maybe Int}.

  \item Quando uma verificação anterior, como \haskellinline{isJust}, já garantiu que o valor não é \haskellinline{Nothing}.


\end{enumerate}

\questionTitle{3}{Falha de conversão com read}
\begin{flushright}
  \footnotesize
  \textit{\textbf{8: Programas interativos}}\\
  \textit{c08-conversao-validacao-02}\\

\end{flushright}
Considere o programa:

\begin{minted}{haskell}
main :: IO ()
main = do
  putStrLn "Digite um número:"
  linha <- getLine
  let numero = read linha :: Int
  print (numero * 2)
\end{minted}

O que acontece se o usuário digitar \haskellinline|"cinco"|?

\begin{enumerate}
  \item O programa chama \haskellinline|getLine| novamente até que o usuário digite um número.

  \item O programa imprime \haskellinline|"cincocinco"|, pois o operador \haskellinline|*| repete strings.

  \item O programa imprime \haskellinline|0|, pois \haskellinline|read| retorna zero quando a conversão falha.

  \item O programa não compila, pois o compilador detecta estaticamente que o usuário pode digitar texto inválido.

  \item A conversão com \haskellinline|read| falha e o programa termina com uma exceção em tempo de execução quando \haskellinline|numero| é avaliado.


\end{enumerate}

\questionTitle{4}{Shadowing em expressões let aninhadas}
\begin{flushright}
  \footnotesize
  \textit{\textbf{3: Expressões e definições}}\\
  \textit{c03-let-escopo-02}\\

\end{flushright}
Considere a expressão:

\begin{haskellcode}
let x = 10
    y = x + 1
in let x = 3
   in x + y
\end{haskellcode}

Qual é o seu valor?

\begin{enumerate}
  \item Ocorre erro de escopo, pois não é permitido reutilizar o nome \haskellinline|x| em um \haskellinline|let| aninhado.

  \item \haskellinline|23|, pois os dois valores de \haskellinline|x| são somados a \haskellinline|y|.

  \item \haskellinline|14|, pois o \haskellinline|x| interno vale \haskellinline|3| no corpo interno, mas \haskellinline|y| foi definido a partir do \haskellinline|x| externo.

  \item \haskellinline|6|, pois o \haskellinline|x| interno redefine também o valor de \haskellinline|y|.

  \item \haskellinline|21|, pois o \haskellinline|x| externo continua valendo em todo o corpo interno.


\end{enumerate}

\questionTitle{5}{Servidor de Linguagem Haskell (HLS)}
\begin{flushright}
  \footnotesize
  \textit{\textbf{2: Ambiente de desenvolvimento}}\\
  \textit{c02-ecossistema-ferramentas-02}\\

\end{flushright}
Ao configurar o editor Visual Studio Code (VS Code) para desenvolvimento Haskell, instala-se o \emph{Haskell Language Server} (HLS). Qual benefício arquitetural essa ferramenta adiciona ao editor?

\begin{enumerate}
  \item Fornece funcionalidades de \emph{IDE} inteligentes, como autocompletar, verificação de erros em tempo real e exibição de tipos ao passar o cursor sobre funções.

  \item Altera as fontes tipográficas e o esquema de cores do editor para padronizar a exibição visual de mônadas e funtores.

  \item Sincroniza os arquivos locais automaticamente com as respostas exigidas no Moodle por meio de \emph{commits} contínuos.

  \item Compila o código-fonte em segundo plano para enviar o binário final diretamente a plataformas de hospedagem em nuvem.

  \item Obriga o VS Code a rodar em um servidor web remoto, permitindo que múltiplos alunos programem no mesmo arquivo simultaneamente.


\end{enumerate}

\questionTitle{6}{Identificação de chamada de cauda}
\begin{flushright}
  \footnotesize
  \textit{\textbf{9: Funções recursivas}}\\
  \textit{c09-cauda-identificacao-01}\\

\end{flushright}
Qual das funções abaixo apresenta a chamada recursiva em posição de cauda?

\begin{enumerate}
  \item \begin{haskellcode}
fib n
  | n == 0 = 0
  | n == 1 = 1
  | n > 1  = fib (n - 2) + fib (n - 1)
\end{haskellcode}

  \item \begin{haskellcode}
produtoIntervalo m n
  | m > n     = 1
  | otherwise = m * produtoIntervalo (m + 1) n
\end{haskellcode}

  \item \begin{haskellcode}
potDois n
  | n == 0 = 1
  | n > 0  = 2 * potDois (n - 1)
\end{haskellcode}

  \item \begin{haskellcode}
fatorial n
  | n == 0 = 1
  | n > 0  = n * fatorial (n - 1)
\end{haskellcode}

  \item \begin{haskellcode}
pdois' x y
  | x == 0 = y
  | x > 0  = pdois' (x - 1) (2 * y)
\end{haskellcode}


\end{enumerate}

\questionTitle{7}{Ambiguidade com read dentro de length}
\begin{flushright}
  \footnotesize
  \textit{\textbf{7: Polimorfismo}}\\
  \textit{cap07-defaulting-length-read}\\

\end{flushright}
Considere:

\begin{minted}{haskell}
r = length [read "10", read "20"]
\end{minted}

Por que essa definição é problemática em compilação normal?

\begin{enumerate}
  \item A definição é sempre resolvida como \haskellinline{[Int]} por causa dos caracteres numéricos nas strings.
  \item \haskellinline{read} sempre retorna \haskellinline{String}, e \haskellinline{length} não aceita listas de strings.
  \item O resultado de \haskellinline{length} é \haskellinline{Int}, mas o tipo dos elementos lidos permanece restrito apenas por \haskellinline{Read a}, sem contexto suficiente para escolher \haskellinline{a}.
  \item O problema é que listas em Haskell não podem conter valores produzidos por funções.
  \item \haskellinline{length} exige que os elementos da lista sejam instâncias de \haskellinline{Num}.

\end{enumerate}

\questionTitle{8}{O uso exploratório do :help}
\begin{flushright}
  \footnotesize
  \textit{\textbf{2: Ambiente de desenvolvimento}}\\
  \textit{c02-ghci-inspecao-02}\\

\end{flushright}
Qual das seguintes opções descreve a função do comando \texttt{:help} (ou \texttt{:?}) no ambiente interativo GHCi?

\begin{enumerate}
  \item Ele formata o arquivo atual para corrigir falhas de indentação antes do carregamento.

  \item Ele aciona o depurador passo a passo (\emph{debugger}) para ajudar a encontrar falhas lógicas em tempo de execução.

  \item Ele exibe um manual interno listando todos os comandos especiais do GHCi e suas respectivas funções.

  \item Ele analisa o código que está sendo digitado e sugere o preenchimento automático das funções em tempo real.

  \item Ele conecta o terminal à internet para buscar soluções de erros de compilação no fórum oficial da linguagem.


\end{enumerate}

\questionTitle{9}{Identificação do caso base em leitura com sentinela}
\begin{flushright}
  \footnotesize
  \textit{\textbf{10: Ações de E/S recursivas}}\\
  \textit{c10-lacos-io-02}\\

\end{flushright}
Considere a ação:

\begin{haskellcode}
lerESomar :: IO Integer
lerESomar = do
  n <- readLn
  if n == 0
    then return 0
    else do
      somaResto <- lerESomar
      return (n + somaResto)
\end{haskellcode}

Qual trecho corresponde ao caso base da recursão?

\begin{enumerate}
  \item \haskellinline|n <- readLn|, pois toda leitura encerra uma iteração.

  \item \haskellinline|return (n + somaResto)|, pois toda expressão com \haskellinline|return| é necessariamente o caso base.

  \item \haskellinline|n == 0| junto com \haskellinline|return 0|, pois o sentinela indica que a leitura deve parar e produz a soma neutra.

  \item O código não possui caso base, pois a ação \haskellinline|lerESomar| aparece no próprio corpo.

  \item \haskellinline|somaResto <- lerESomar|, pois a chamada recursiva define quando a repetição termina.


\end{enumerate}

\questionTitle{10}{Identificando tipos vs. variáveis}
\begin{flushright}
  \footnotesize
  \textit{\textbf{4: Tipos de Dados}}\\
  \textit{c04-conceito-nomes-02}\\

\end{flushright}
Considere os identificadores: \haskellinline{Pessoa} e \haskellinline{pessoa}. Qual é a interpretação mais provável em Haskell?

\begin{enumerate}
  \item Ambos são tipos, sendo \haskellinline{pessoa} um sinônimo.

  \item Ambos são variáveis, mas com escopos diferentes.

  \item \haskellinline{Pessoa} é uma variável global e \haskellinline{pessoa} é uma variável local.

  \item \haskellinline{Pessoa} é um tipo (ou construtor de dados) e \haskellinline{pessoa} é uma variável ou função.

  \item Não há diferença; Haskell não distingue maiúsculas de minúsculas.


\end{enumerate}

\questionTitle{11}{Quantidade de leituras até o sentinela}
\begin{flushright}
  \footnotesize
  \textit{\textbf{10: Ações de E/S recursivas}}\\
  \textit{c10-rastreamento-02}\\

\end{flushright}
Considere a mesma função:

\begin{haskellcode}
lerESomar :: IO Integer
lerESomar = do
  n <- readLn
  if n == 0
    then return 0
    else do
      somaResto <- lerESomar
      return (n + somaResto)
\end{haskellcode}

Se o usuário digita \haskellinline|2, 2, 0|, quantas chamadas a \haskellinline|readLn| são realizadas?

\begin{enumerate}
  \item 0

  \item 2

  \item 3

  \item 1

  \item 4


\end{enumerate}

\questionTitle{12}{Garantias da função pura}
\begin{flushright}
  \footnotesize
  \textit{\textbf{1: Paradigmas de programação}}\\
  \textit{c01-transparencia-referencial-03}\\

\end{flushright}
Uma \textbf{função pura} é a base fundamental da construção de software no paradigma declarativo funcional. Assinale a alternativa que descreve corretamente as duas garantias absolutas oferecidas por uma função pura.

\begin{enumerate}
  \item 1) Ela não utiliza declarações de variáveis locais em seu corpo. 2) Ela só pode receber matrizes e listas imutáveis como parâmetros.

  \item 1) Ela sempre avalia para o mesmo resultado quando recebe os mesmos argumentos. 2) Ela não causa nenhum efeito colateral observável fora de seu escopo.

  \item 1) Ela avalia todos os seus argumentos de forma estrita (\emph{strict evaluation}) antes de executar. 2) Ela não depende de bibliotecas de terceiros.

  \item 1) Ela é garantida de terminar (não entrar em loop infinito). 2) Ela executa obrigatoriamente em tempo constante $\mathcal{O}(1)$.

  \item 1) Ela possui uma assinatura de tipo polimórfica explícita. 2) Ela substitui automaticamente operações de ponto flutuante por inteiros para evitar imprecisões matemáticas.


\end{enumerate}

\questionTitle{13}{Chamada que se afasta do caso base}
\begin{flushright}
  \footnotesize
  \textit{\textbf{9: Funções recursivas}}\\
  \textit{c09-fundamentos-terminacao-02}\\

\end{flushright}
Considere a função:
\begin{haskellcode}
potDois' :: Integer -> Integer
potDois' n
  | n == 0    = 1
  | otherwise = 2 * potDois' (n - 1)
\end{haskellcode}
Qual é o problema ao avaliar \haskellinline|potDois' (-5)|?

\begin{enumerate}
  \item A função retorna \haskellinline|1|, pois todo expoente negativo é tratado como o caso base.

  \item A função retorna \haskellinline|-32|, pois Haskell preserva automaticamente o sinal do expoente.

  \item Ocorre erro de tipo, pois \haskellinline|Integer| não permite valores negativos.

  \item A função retorna \haskellinline|0|, pois potências de dois negativas são arredondadas para zero.

  \item A cada chamada o argumento fica menor, afastando-se de \haskellinline|0|; portanto, a recursão não termina.


\end{enumerate}

\questionTitle{14}{Tipo produzido por return no caso base}
\begin{flushright}
  \footnotesize
  \textit{\textbf{10: Ações de E/S recursivas}}\\
  \textit{c10-return-tipos-02}\\

\end{flushright}
Analise a ação recursiva:

\begin{haskellcode}
lerLista :: IO [Int]
lerLista = do
  x <- readLn
  if x == 0
    then return []
    else do
      resto <- lerLista
      return (x:resto)
\end{haskellcode}

Tecnicamente, o que a expressão \haskellinline|return []| realiza no caso base?

\begin{enumerate}
  \item Ela lê a próxima lista da entrada padrão, usando \haskellinline|readLn| implicitamente.

  \item Ela sinaliza ao compilador que a recursão deve ser transformada em um laço imperativo.

  \item Ela imprime \haskellinline|[]| na tela e retorna \haskellinline|()|.

  \item Ela encerra todo o programa e descarta as chamadas recursivas anteriores.

  \item Ela cria uma ação de tipo \haskellinline|IO [Int]| que produz a lista vazia quando executada.


\end{enumerate}

\questionTitle{15}{Restrição Num e ausência de coerção implícita}
\begin{flushright}
  \footnotesize
  \textit{\textbf{7: Polimorfismo}}\\
  \textit{cap07-adhoc-restricao-num}\\

\end{flushright}
Considere a assinatura:

\begin{minted}{haskell}
(+) :: Num a => a -> a -> a
\end{minted}

Qual interpretação é a mais precisa?

\begin{enumerate}
  \item A assinatura diz que \haskellinline{(+)} é paramétrica pura, sem qualquer restrição sobre \haskellinline{a}.
  \item Os dois operandos e o resultado têm o mesmo tipo \haskellinline{a}; esse tipo deve ser instância de \haskellinline{Num}.
  \item O compilador escolhe uma implementação por herança de classe de POO.
  \item A restrição \haskellinline{Num a} significa que o resultado será sempre \haskellinline{Integer}.
  \item O operador aceita livremente um \haskellinline{Int} e um \haskellinline{Double}, pois ambos pertencem à hierarquia numérica.

\end{enumerate}

\questionTitle{16}{Conflito de classes em checagem estática}
\begin{flushright}
  \footnotesize
  \textit{\textbf{4: Tipos de Dados}}\\
  \textit{c04-checagem-conflito-01}\\

\end{flushright}
Considere o código \haskellinline{calculo x y = (x / 2) `div` y}. Por que este código falha na compilação?

\begin{enumerate}
  \item O operador \haskellinline{/} exige um tipo fracionário (\haskellinline{Fractional}), mas \haskellinline{div} exige um tipo integral (\haskellinline{Integral}), causando um conflito de tipos.

  \item A função \haskellinline{div} tem maior precedência que \haskellinline{/}.

  \item O literal \haskellinline{2} é inferido como \haskellinline{Integer}, que não é compatível com \haskellinline{x}.

  \item Falta uma anotação de tipo explícita, impedindo a inferência.

  \item A notação infixa com crases não pode ser usada com operadores simbólicos.


\end{enumerate}

\questionTitle{17}{Carregando módulos no GHCi}
\begin{flushright}
  \footnotesize
  \textit{\textbf{2: Ambiente de desenvolvimento}}\\
  \textit{c02-ghci-modulos-02}\\

\end{flushright}
Para utilizar funções avançadas de listas (como \haskellinline{sort}) diretamente no prompt do GHCi, o desenvolvedor precisa trazer o módulo \texttt{Data.List} para o escopo atual. Quais comandos são válidos para realizar esta operação?

\begin{enumerate}
  \item Não é necessário usar comandos; em Haskell, diferentemente de C ou Java, todos os módulos oficiais estão sempre injetados globalmente no Prelude.

  \item Apenas o comando interno \texttt{:load Data.List}, pois o termo \haskellinline{import} é exclusivo para arquivos físicos \texttt{.hs}.

  \item O comando \texttt{:browse Data.List}, que carrega e executa todas as funções contidas no módulo de forma simultânea.

  \item A palavra-chave nativa da linguagem \haskellinline{import Data.List} ou o atalho específico do GHCi \texttt{:module +Data.List} (ou \texttt{:m +Data.List}).

  \item O comando \texttt{:install Data.List}, que solicita a função diretamente do repositório remoto Hackage em tempo real.


\end{enumerate}

\questionTitle{18}{Declaração let no GHCi}
\begin{flushright}
  \footnotesize
  \textit{\textbf{3: Expressões e definições}}\\
  \textit{c03-ghci-definicoes-02}\\

\end{flushright}
No GHCi, qual alternativa representa a forma explícita de associar o nome \haskellinline|raio| ao valor \haskellinline|10| para uso em expressões posteriores?

\begin{enumerate}
  \item \haskellinline|let 10 = raio|.

  \item \haskellinline|raio := 10|.

  \item \haskellinline|let raio = 10|.

  \item \haskellinline|def raio = 10|.

  \item \haskellinline|raio <- 10|.


\end{enumerate}

\questionTitle{19}{Análise de IMC com where}
\begin{flushright}
  \footnotesize
  \textit{\textbf{5: Expressão condicional}}\\
  \textit{c05-where-avaliacao-03}\\

\end{flushright}
Considere a função:
\begin{haskellcode}
analisaIMC peso altura
  | imc < 18.5 = "Abaixo"
  | imc < 25.0 = "Normal"
  | otherwise  = "Acima"
  where
    imc = peso / (altura * altura)
\end{haskellcode}
Qual é o resultado de \haskellinline|analisaIMC 80 2.0|?

\begin{enumerate}
  \item \haskellinline|"80.0"|

  \item \haskellinline|"Abaixo"|

  \item Ocorre um erro, pois variáveis definidas em \haskellinline|where| não podem aparecer nas guardas.

  \item \haskellinline|"Acima"|

  \item \haskellinline|"Normal"|


\end{enumerate}

\questionTitle{20}{Inferência de tipo de função composta}
\begin{flushright}
  \footnotesize
  \textit{\textbf{4: Tipos de Dados}}\\
  \textit{c04-inferencia-funcao-01}\\

\end{flushright}
Considere a função definida por \haskellinline{f x = not (x > 0)}. Qual é o tipo inferido para \haskellinline{f}?

\begin{enumerate}
  \item \haskellinline{Double -> Bool}

  \item \haskellinline{(Ord a, Num a) => a -> Bool}

  \item \haskellinline{Int -> Bool}

  \item \haskellinline{Num a => a -> Bool}

  \item \haskellinline{a -> Bool}


\end{enumerate}

\questionTitle{21}{If produzindo função}
\begin{flushright}
  \footnotesize
  \textit{\textbf{5: Expressão condicional}}\\
  \textit{c05-if-avaliacao-03}\\

\end{flushright}
Qual é o resultado da expressão abaixo?
\begin{haskellcode}
(if 5 > 3 then (*) else (+)) 4 2
\end{haskellcode}

\begin{enumerate}
  \item \haskellinline|6|

  \item Ocorre um erro de tipo, pois os ramos do \haskellinline|if| são funções.

  \item \haskellinline|8|

  \item Ocorre um erro de sintaxe, pois \haskellinline|if| não pode aparecer entre parênteses.

  \item \haskellinline|14|


\end{enumerate}

\questionTitle{22}{Tipos de ações de entrada}
\begin{flushright}
  \footnotesize
  \textit{\textbf{8: Programas interativos}}\\
  \textit{c08-tipos-acoes-02}\\

\end{flushright}
Qual é a assinatura de tipo correta de \haskellinline|getLine| e o que ela representa?

\begin{enumerate}
  \item \haskellinline|IO String| — uma ação de E/S que, quando executada, lê uma linha de texto da entrada padrão.

  \item \haskellinline|IO ()| — uma ação que executa a entrada, mas descarta qualquer valor produzido.

  \item \haskellinline|String -> String| — uma função pura que recebe uma string e retorna o valor lido.

  \item \haskellinline|() -> IO String| — uma função que recebe o valor unitário e então constrói a ação.

  \item \haskellinline|String| — um valor puro obtido diretamente, sem necessidade de \haskellinline|IO|.


\end{enumerate}

\questionTitle{23}{Estratégias de Redução e Confluência}
\begin{flushright}
  \footnotesize
  \textit{\textbf{1: Paradigmas de programação}}\\
  \textit{c01-mecanismo-reducao-03}\\

\end{flushright}
Ao avaliar uma expressão matemática em Haskell, o sistema pode teoricamente adotar diferentes caminhos de avaliação, como a redução \emph{de dentro para fora} (aplicativa) ou \emph{de fora para dentro} (normal/preguiçosa). 

Graças à pureza e à ausência de efeitos colaterais na linguagem, qual garantia teórica fundamental o programador possui independentemente do caminho escolhido pelo compilador?

\begin{enumerate}
  \item A garantia de que os tipos numéricos fracionários estarão imunes a qualquer tipo de erro de precisão de ponto flutuante do padrão IEEE 754.

  \item A garantia de que a quantidade de memória RAM alocada dinamicamente durante o processo de avaliação será anulada pelo \emph{Garbage Collector}.

  \item A garantia de que, caso ambas as estratégias cheguem a um resultado final sem entrar em \emph{loop} infinito, este resultado será sempre o mesmo (Teorema de Church-Rosser).

  \item A garantia de que as funções envolvidas na expressão serão automaticamente vetorizadas e enviadas para execução na GPU.

  \item A garantia de que a execução da expressão consumirá exatamente a mesma quantidade de milissegundos na CPU do usuário.


\end{enumerate}

\questionTitle{24}{O papel de hFlush}
\begin{flushright}
  \footnotesize
  \textit{\textbf{8: Programas interativos}}\\
  \textit{c08-buffering-02}\\

\end{flushright}
Qual ação pode ser usada imediatamente depois de \haskellinline|putStr| para garantir que o conteúdo pendente na saída padrão seja enviado ao terminal?

\begin{enumerate}
  \item \haskellinline|getLine stdout|.

  \item \haskellinline|hRefresh stdout|.

  \item \haskellinline|delay stdout|.

  \item \haskellinline|hClear stdout|.

  \item \haskellinline|hFlush stdout|.


\end{enumerate}

\questionTitle{25}{Precedência da aplicação sobre operadores}
\begin{flushright}
  \footnotesize
  \textit{\textbf{3: Expressões e definições}}\\
  \textit{c03-precedencia-associatividade-02}\\

\end{flushright}
A expressão \haskellinline|sqrt 16 + 3 * 2| ilustra a precedência da aplicação de função em relação aos operadores infixos. Qual alternativa descreve corretamente sua avaliação?

\begin{enumerate}
  \item Ela é lida como \haskellinline|sqrt (16 + 3 * 2)|, resultando em aproximadamente \haskellinline|4.6904|.

  \item Ela é lida como \haskellinline|(sqrt 16) + (3 * 2)|, resultando em \haskellinline|10.0|.

  \item Ela é inválida porque operadores aritméticos não podem ser misturados com aplicação prefixa.

  \item Ela é inválida porque \haskellinline|sqrt| exige parênteses em volta de \haskellinline|16|.

  \item Ela é lida como \haskellinline|sqrt (16 + 3) * 2|, resultando em aproximadamente \haskellinline|8.7178|.


\end{enumerate}

\questionTitle{26}{Rastreamento da potência de dois}
\begin{flushright}
  \footnotesize
  \textit{\textbf{9: Funções recursivas}}\\
  \textit{c09-rastreamento-potdois-01}\\

\end{flushright}
Considere:
\begin{haskellcode}
potDois :: Integer -> Integer
potDois n
  | n == 0 = 1
  | n > 0  = 2 * potDois (n - 1)
\end{haskellcode}
Qual é o valor de \haskellinline|potDois 4|?

\begin{enumerate}
  \item \haskellinline|24|

  \item \haskellinline|12|

  \item \haskellinline|8|

  \item \haskellinline|32|

  \item \haskellinline|16|


\end{enumerate}

\questionTitle{27}{Papel de where em guardas}
\begin{flushright}
  \footnotesize
  \textit{\textbf{5: Expressão condicional}}\\
  \textit{c05-where-escopo-03}\\

\end{flushright}
Qual é a principal vantagem de usar \haskellinline|where| em uma definição com guardas, como no cálculo de \haskellinline|delta| em uma função que determina o número de raízes de uma equação?

\begin{enumerate}
  \item Tornar a função automaticamente polimórfica.

  \item Permitir nomear um cálculo local e reutilizá-lo em guardas e resultados, evitando repetição.

  \item Fazer com que as guardas sejam avaliadas em paralelo.

  \item Permitir que a função tenha ramos com tipos diferentes.

  \item Garantir que todas as guardas sejam exaustivas.


\end{enumerate}

\questionTitle{28}{A inovação do Sistema de Tipos da ML}
\begin{flushright}
  \footnotesize
  \textit{\textbf{1: Paradigmas de programação}}\\
  \textit{c01-historia-evolucao-03}\\

\end{flushright}
Na década de 1970, o pesquisador Robin Milner e sua equipe desenvolveram a linguagem ML (\emph{MetaLanguage}), projetada inicialmente para apoiar a demonstração automática de teoremas. A linguagem ML legou um avanço formidável, hoje presente fortemente em Haskell e linguagens como Rust e Swift. Qual foi essa inovação técnica?

\begin{enumerate}
  \item O mecanismo de \textbf{Avaliação Preguiçosa} estrita, que impedia o cálculo precipitado de matrizes infinitas e protegia os \emph{mainframes} contra falhas de estouro de memória.

  \item Um sistema de \textbf{inferência de tipos} elegante e poderoso, capaz de garantir a segurança estática sem exigir que o programador preenchesse o código com anotações manuais redundantes de tipos.

  \item A invenção pioneira da \textbf{Orientação a Objetos}, que permitiu à indústria mesclar funções puras com o encapsulamento hermético de estado dentro de classes herdáveis.

  \item A introdução e formalização das \textbf{Mônadas} na biblioteca padrão como a única abstração autorizada pelo compilador para isolar e manipular efeitos de entrada e saída.

  \item A arquitetura inovadora do compilador \textbf{JIT} (\emph{Just-in-Time Compiler}), que otimizou o código funcional para superar a velocidade dos programas escritos em linguagem C de forma absoluta.


\end{enumerate}

\questionTitle{29}{Inferência de literais no GHCi (parametrizado)}
\begin{flushright}
  \footnotesize
  \textit{\textbf{4: Tipos de Dados}}\\
  \textit{c04-ghci-inferencia-01}\\

\end{flushright}
Qual tipo o GHCi inferiria para o literal \haskellinline{True } ao usar o comando \texttt{:type}?

\begin{enumerate}
  \item \haskellinline{Bit}

  \item \haskellinline{Bool}

  \item \haskellinline{Int}

  \item Ocorre um erro de sintaxe.

  \item \haskellinline{Boolean}


\end{enumerate}

\questionTitle{30}{Instanciação de variáveis de tipo em map}
\begin{flushright}
  \footnotesize
  \textit{\textbf{7: Polimorfismo}}\\
  \textit{cap07-param-map-instanciacao}\\

\end{flushright}
Considere a assinatura:

\begin{minted}{haskell}
map :: (a -> b) -> [a] -> [b]
\end{minted}

Na expressão abaixo, quais instanciações de \haskellinline{a} e \haskellinline{b} são usadas?

\begin{minted}{haskell}
map even [1, 2, 3, 4]
\end{minted}

\begin{enumerate}
  \item \haskellinline{a} é instanciado como um tipo integral, usualmente \haskellinline{Integer} no contexto padrão do GHCi, e \haskellinline{b} é instanciado como \haskellinline{Bool}.
  \item \haskellinline{a} é instanciado como \haskellinline{Bool} e \haskellinline{b} como \haskellinline{Integer}, pois \haskellinline{even} retorna números pares.
  \item A expressão não é tipável, pois \haskellinline{map} só aceita listas de funções.
  \item \haskellinline{a} e \haskellinline{b} são ambos instanciados como \haskellinline{Bool}, pois o resultado de \haskellinline{even} é booleano.
  \item \haskellinline{a} e \haskellinline{b} permanecem totalmente independentes, pois \haskellinline{map} nunca usa o tipo da função recebida.

\end{enumerate}

\questionTitle{31}{Tipo de uma função com igualdade e literal}
\begin{flushright}
  \footnotesize
  \textit{\textbf{7: Polimorfismo}}\\
  \textit{cap07-expressao-eq-zero}\\

\end{flushright}
Qual é a assinatura mais geral da função abaixo?

\begin{minted}{haskell}
zeros xs = map (== 0) xs
\end{minted}

\begin{enumerate}
  \item \mintinline{haskell}{zeros :: Eq a => [a] -> [Bool]}
  \item \mintinline{haskell}{zeros :: Num a => [a] -> [a]}
  \item \mintinline{haskell}{zeros :: [Bool] -> Bool}
  \item \mintinline{haskell}{zeros :: [Int] -> [Int]}
  \item \mintinline{haskell}{zeros :: (Eq a, Num a) => [a] -> [Bool]}

\end{enumerate}

\questionTitle{32}{print e putStrLn}
\begin{flushright}
  \footnotesize
  \textit{\textbf{8: Programas interativos}}\\
  \textit{c08-entrada-saida-02}\\

\end{flushright}
Qual é a diferença central entre \haskellinline|print| e \haskellinline|putStrLn|?

\begin{enumerate}
  \item \haskellinline|print| lê um valor da entrada; \haskellinline|putStrLn| imprime um valor na saída.

  \item \haskellinline|print| aplica \haskellinline|show| ao valor antes de imprimir; \haskellinline|putStrLn| recebe diretamente uma \haskellinline|String| e acrescenta uma quebra de linha.

  \item Não há diferença; as duas funções têm exatamente o mesmo tipo e o mesmo comportamento.

  \item \haskellinline|print| imprime sem quebra de linha; \haskellinline|putStrLn| imprime com quebra de linha.

  \item \haskellinline|print| só funciona para \haskellinline|String|; \haskellinline|putStrLn| funciona para qualquer tipo da classe \haskellinline|Show|.


\end{enumerate}

\questionTitle{33}{O papel de main}
\begin{flushright}
  \footnotesize
  \textit{\textbf{8: Programas interativos}}\\
  \textit{c08-modelo-io-02}\\

\end{flushright}
Qual é o papel especial de uma definição como \haskellinline|main :: IO ()| em um programa Haskell executável?

\begin{enumerate}
  \item \haskellinline|main| é apenas uma convenção de nomenclatura; qualquer função do arquivo poderia ser executada automaticamente.

  \item \haskellinline|main| é a única definição do programa que pode chamar funções puras.

  \item \haskellinline|main| é obrigatoriamente uma função pura, pois programas Haskell não podem ter efeitos de E/S.

  \item \haskellinline|main| é a ação de E/S que o \emph{runtime system} executa quando o programa é iniciado.

  \item \haskellinline|main| deve retornar um código inteiro de saída, como ocorre em C.


\end{enumerate}

\questionTitle{34}{Paradigmas - Princípio Nuclear (Legado Refatorado)}
\begin{flushright}
  \footnotesize
  \textit{\textbf{1: Paradigmas de programação}}\\
  \textit{c01-conceitos-paradigmas-02}\\

\end{flushright}
Ao contrastarmos os paradigmas dominantes, o \textbf{princípio nuclear} que separa a programação imperativa da programação declarativa (funcional) reside no foco da estruturação do código. Qual é essa diferença?

\begin{enumerate}
  \item A programação imperativa é estruturada em torno de funções matemáticas puras, enquanto a declarativa é estruturada em torno de objetos que trocam mensagens.

  \item A programação imperativa foca em descrever \emph{como} o computador deve executar os passos (controle de fluxo e estado), enquanto a declarativa foca em descrever \emph{o que} é o resultado desejado (lógica e propriedades).

  \item A programação imperativa abstrai o gerenciamento de memória através de \emph{Garbage Collectors}, enquanto a declarativa exige alocação explícita com \cinline|malloc|.

  \item A programação imperativa exige a declaração estrita de tipos em tempo de compilação, enquanto a declarativa permite que os tipos sejam inferidos dinamicamente durante a execução.

  \item A programação imperativa foi projetada exclusivamente para arquiteturas mononucleares (single-core), enquanto a declarativa foi projetada nativamente para paralelismo em GPUs.


\end{enumerate}

\questionTitle{35}{Por que Int não satisfaz Fractional}
\begin{flushright}
  \footnotesize
  \textit{\textbf{7: Polimorfismo}}\\
  \textit{cap07-hierarquia-int-fracao}\\

\end{flushright}
Por que a expressão abaixo não é aceita sem conversão explícita?

\begin{minted}{haskell}
n :: Int
n = 10

n / 2
\end{minted}

\begin{enumerate}
  \item Porque Haskell converte automaticamente \haskellinline{Int} para \haskellinline{Double}, mas apenas fora do GHCi.
  \item Porque \haskellinline{(/)} exige \haskellinline{Fractional a => a -> a -> a}, e \haskellinline{Int} não é instância de \haskellinline{Fractional}.
  \item Porque \haskellinline{Int} não é instância de \haskellinline{Num}.
  \item Porque literais inteiros nunca podem aparecer em expressões com divisão.
  \item Porque \haskellinline{(/)} só existe para \haskellinline{Integer}.

\end{enumerate}

\questionTitle{36}{Ausência de coerção implícita entre Int e Double}
\begin{flushright}
  \footnotesize
  \textit{\textbf{7: Polimorfismo}}\\
  \textit{cap07-conversao-int-double-soma}\\

\end{flushright}
Considere:

\begin{minted}{haskell}
n :: Int
n = 10

x :: Double
x = 2.5
\end{minted}

Qual expressão soma corretamente os valores como \haskellinline{Double}?

\begin{enumerate}
  \item \mintinline{haskell}{read n + x}
  \item \mintinline{haskell}{n + x}
  \item \mintinline{haskell}{show n + x}
  \item \mintinline{haskell}{fromIntegral n + x}
  \item \mintinline{haskell}{toInteger x + n}

\end{enumerate}

\questionTitle{37}{Próxima equação da função}
\begin{flushright}
  \footnotesize
  \textit{\textbf{5: Expressão condicional}}\\
  \textit{c05-guardas-exaustividade-03}\\

\end{flushright}
Considere a definição:
\begin{haskellcode}
h x
  | x > 0 = "Positivo"
h x = "Não positivo"
\end{haskellcode}
Qual é o resultado de \haskellinline|h 0|?

\begin{enumerate}
  \item O compilador rejeita a definição por possuir duas equações para a mesma função.

  \item Ocorre um erro, pois a primeira equação não tem \haskellinline|otherwise|.

  \item \haskellinline|0|

  \item \haskellinline|"Não positivo"|

  \item \haskellinline|"Positivo"|


\end{enumerate}

\questionTitle{38}{Uso do tipo Rational}
\begin{flushright}
  \footnotesize
  \textit{\textbf{4: Tipos de Dados}}\\
  \textit{c04-tipos-rational-01}\\

\end{flushright}
O tipo \haskellinline{Rational} representa números racionais como uma fração de dois \haskellinline{Integer}. Como se constrói um valor do tipo \haskellinline{Rational} em Haskell?

\begin{enumerate}
  \item Usando o operador \haskellinline{/}, por exemplo: \haskellinline{3 / 4}.

  \item \haskellinline{Rational} é apenas um sinônimo para \haskellinline{Double}.

  \item Usando a função \haskellinline{ratio 3 4} do Prelude.

  \item Usando o operador \haskellinline{%} do módulo \haskellinline{Data.Ratio}, por exemplo: \haskellinline{3 % 4}.

  \item Escrevendo o literal \haskellinline{3/4} entre aspas.


\end{enumerate}

\questionTitle{39}{Paralelismo Seguro}
\begin{flushright}
  \footnotesize
  \textit{\textbf{1: Paradigmas de programação}}\\
  \textit{c01-engenharia-software-03}\\

\end{flushright}
No contexto de sistemas de alta disponibilidade executados em processadores \emph{multicore}, a programação funcional tem ganhado vasto espaço na indústria contemporânea. Qual característica fundacional do paradigma facilita imensamente a construção de programas paralelos livres de falhas?

\begin{enumerate}
  \item A imutabilidade dos dados combinada às funções puras garantem que múltiplas \emph{threads} não causarão \emph{race conditions} (condições de corrida) tentando modificar o mesmo espaço de memória simultaneamente.

  \item A tipagem dinâmica fraca da linguagem assegura que os processos em segundo plano reescrevam as estruturas de memória de forma atômica e sincronizada por meio do \emph{garbage collector}.

  \item O paradigma impossibilita deliberadamente a criação de processos paralelos (sendo restrito a \emph{single-thread}), o que mitiga os erros de concorrência anulando a complexidade do escalonamento.

  \item Os compiladores de linguagens funcionais traduzem nativamente o código para chamadas assíncronas de \emph{hardware} que contornam totalmente os barramentos lentos da memória principal.

  \item A avaliação estrita das linguagens funcionais impõe um controle rígido ao escalonador do sistema operacional, forçando a distribuição milimétrica da carga da aplicação entre os núcleos físicos.


\end{enumerate}

\questionTitle{40}{Modelagem composta de pedidos}
\begin{flushright}
  \footnotesize
  \textit{\textbf{6: Estruturas de dados básicas}}\\
  \textit{c06-modelagem-03}\\

\end{flushright}
Qual tipo representa melhor uma coleção de pedidos em que cada pedido possui o nome do cliente, a lista de itens comprados e um total opcional?

\begin{enumerate}
  \item \haskellinline{[String, [(String, Int)], Maybe Double]}

  \item \haskellinline{(String, [String], Double)}

  \item \haskellinline{Maybe [(String, Int, Double)]}

  \item \haskellinline{[(String, [(String, Int)], Maybe Double)]}

  \item \haskellinline{([String], [Int], [Double])}


\end{enumerate}

\questionTitle{41}{Comentários aninhados e código efetivo}
\begin{flushright}
  \footnotesize
  \textit{\textbf{3: Expressões e definições}}\\
  \textit{c03-comentarios-valores-02}\\

\end{flushright}
Considere o trecho a seguir.

\begin{haskellcode}
x = 10
{-
y = 20
  {- z = 30 -}
w = 40
-}
r = x + 1
\end{haskellcode}

Quais definições permanecem como código efetivo para o compilador?

\begin{enumerate}
  \item Nenhuma definição, pois comentários de bloco aninhados invalidam o arquivo.

  \item \haskellinline|x|, \haskellinline|z| e \haskellinline|r|.

  \item Apenas \haskellinline|r|, pois \haskellinline|x| fica dentro do comentário de bloco.

  \item Apenas \haskellinline|x| e \haskellinline|r|.

  \item \haskellinline|x|, \haskellinline|y|, \haskellinline|w| e \haskellinline|r|.


\end{enumerate}

\questionTitle{42}{Cons versus concatenação}
\begin{flushright}
  \footnotesize
  \textit{\textbf{6: Estruturas de dados básicas}}\\
  \textit{c06-listas-estrutura-03}\\

\end{flushright}
Qual afirmação descreve corretamente a diferença entre \haskellinline{(:)} e \haskellinline{(++)}?

\begin{enumerate}
  \item \haskellinline{(++)} adiciona um elemento à lista, enquanto \haskellinline{(:)} exige duas listas como argumentos.

  \item \haskellinline{(:)} e \haskellinline{(++)} sempre produzem exatamente o mesmo efeito, mudando apenas a sintaxe.

  \item \haskellinline{(:)} adiciona um elemento ao início de uma lista, enquanto \haskellinline{(++)} concatena duas listas.

  \item \haskellinline{(:)} adiciona um elemento ao final da lista, enquanto \haskellinline{(++)} adiciona ao início.

  \item \haskellinline{(:)} concatena listas mais rapidamente do que \haskellinline{(++)}, mas exige listas do mesmo tamanho.


\end{enumerate}

\questionTitle{43}{Produto de intervalo vazio}
\begin{flushright}
  \footnotesize
  \textit{\textbf{9: Funções recursivas}}\\
  \textit{c09-projeto-produto-intervalo-01}\\

\end{flushright}
Na definição de \haskellinline|produtoIntervalo m n|, adotou-se a convenção de que o produto de um intervalo vazio é \haskellinline|1|. Qual caso base implementa essa convenção?

\begin{enumerate}
  \item \haskellinline|m == 0 = n|

  \item \haskellinline|m == n = 1|

  \item \haskellinline|m < n = 1|

  \item \haskellinline|n == 0 = 0|

  \item \haskellinline|m > n = 1|


\end{enumerate}

\questionTitle{44}{Enumeração com passo e consumo finito}
\begin{flushright}
  \footnotesize
  \textit{\textbf{6: Estruturas de dados básicas}}\\
  \textit{c06-listas-operacoes-03}\\

\end{flushright}
Considere a expressão:
\begin{minted}{haskell}
take 5 [5, 5 + 2 ..]
\end{minted}
Qual é o seu resultado?

\begin{enumerate}
  \item \haskellinline{ [5,10,15,20,25] }

  \item A expressão produz a lista infinita completa antes de aplicar \haskellinline{take}.

  \item \haskellinline{ [5,7,9,11,13] }

  \item \haskellinline{ [7,9,11,13,15] }

  \item \haskellinline{ [5,7,9,11] }


\end{enumerate}

\questionTitle{45}{Assinatura de função com Bool e lista de Char}
\begin{flushright}
  \footnotesize
  \textit{\textbf{4: Tipos de Dados}}\\
  \textit{c04-funcao-assinatura-02}\\

\end{flushright}
Qual das seguintes assinaturas de tipo representa uma função que recebe um \haskellinline{Bool} e uma lista de \haskellinline{Char} e retorna um \haskellinline{Char}?

\begin{enumerate}
  \item \haskellinline{Bool -> [Char] -> Char}

  \item \haskellinline{(Bool, [Char]) -> Char}

  \item \haskellinline{(Bool, [Char], Char)}

  \item \haskellinline{[Bool] -> [Char] -> Char}

  \item \haskellinline{Char -> Bool -> [Char]}


\end{enumerate}

\questionTitle{46}{Erro comum com let}
\begin{flushright}
  \footnotesize
  \textit{\textbf{8: Programas interativos}}\\
  \textit{c08-blocos-do-02}\\

\end{flushright}
Um programador escreveu:

\begin{minted}{haskell}
main :: IO ()
main = do
  putStrLn "Qual é o seu nome?"
  let nome = getLine
  putStrLn ("Olá, " ++ nome)
\end{minted}

Qual é a causa fundamental do erro?

\begin{enumerate}
  \item \haskellinline|nome| fica associado à ação \haskellinline|getLine|, de tipo \haskellinline|IO String|. O operador \haskellinline|++| espera uma \haskellinline|String|, não uma ação \haskellinline|IO String|.

  \item \haskellinline|getLine| não pode ser usado depois de \haskellinline|putStrLn| no mesmo bloco \haskellinline|do|.

  \item A função \haskellinline|putStrLn| não aceita expressões com concatenação de strings.

  \item A variável \haskellinline|nome| precisava ser declarada globalmente antes de \haskellinline|main|.

  \item A palavra \haskellinline|let| é sempre proibida dentro de blocos \haskellinline|do|.


\end{enumerate}

\questionTitle{47}{Inferência com comparação e conversão para String}
\begin{flushright}
  \footnotesize
  \textit{\textbf{7: Polimorfismo}}\\
  \textit{cap07-inferencia-restricoes-ord-show}\\

\end{flushright}
Considere:

\begin{minted}{haskell}
diagnostico x y =
  if x <= y then show x else show y
\end{minted}

Qual é a assinatura mais geral de \haskellinline{diagnostico}?

\begin{enumerate}
  \item \mintinline{haskell}{diagnostico :: Eq a => a -> a -> Bool}
  \item \mintinline{haskell}{diagnostico :: (Ord a, Show a) => a -> a -> String}
  \item \mintinline{haskell}{diagnostico :: Show a => a -> a -> String}
  \item \mintinline{haskell}{diagnostico :: a -> a -> String}
  \item \mintinline{haskell}{diagnostico :: Ord a => a -> a -> String}

\end{enumerate}

\questionTitle{48}{Vantagem de tipos sinônimos}
\begin{flushright}
  \footnotesize
  \textit{\textbf{4: Tipos de Dados}}\\
  \textit{c04-sinonimos-vantagem-01}\\

\end{flushright}
Qual é o principal benefício de usar um tipo sinônimo, como \haskellinline{type Idade = Int}?

\begin{enumerate}
  \item Criar um tipo completamente novo que é incompatível com \haskellinline{Int}, aumentando a segurança.

  \item Aumentar a legibilidade e a expressividade do código, tornando as assinaturas mais auto-documentadas.

  \item Melhorar o desempenho do programa, pois sinônimos são mais eficientes.

  \item Permitir a definição de tipos recursivos.

  \item É a única maneira de usar \haskellinline{Int} em assinaturas de função.


\end{enumerate}

\questionTitle{49}{Guardas sobrepostas}
\begin{flushright}
  \footnotesize
  \textit{\textbf{5: Expressão condicional}}\\
  \textit{c05-guardas-sintaxe-semantica-03}\\

\end{flushright}
Considere a função:
\begin{haskellcode}
f x
  | x > 0     = 1
  | x > 10    = 2
  | otherwise = 3
\end{haskellcode}
Qual é o valor de \haskellinline|f 20|?

\begin{enumerate}
  \item Ocorre um erro, pois duas guardas são verdadeiras.

  \item \haskellinline|2|

  \item \haskellinline|3|

  \item \haskellinline|1|

  \item O compilador rejeita a função, pois as guardas se sobrepõem.


\end{enumerate}

\questionTitle{50}{Soma de dois números}
\begin{flushright}
  \footnotesize
  \textit{\textbf{8: Programas interativos}}\\
  \textit{c08-programas-completos-02}\\

\end{flushright}
Considere:

\begin{minted}{haskell}
main :: IO ()
main = do
  x <- readLn :: IO Int
  y <- readLn :: IO Int
  print (x + y)
\end{minted}

Se o usuário digitar \haskellinline|3| e depois \haskellinline|4|, qual será a saída produzida por \haskellinline|print|?

\begin{enumerate}
  \item \haskellinline|"34"|

  \item Nenhuma saída, pois \haskellinline|readLn| não pode ser usado dentro de \haskellinline|do|.

  \item \haskellinline|(3,4)|

  \item \haskellinline|7|

  \item \haskellinline|3 + 4|


\end{enumerate}

\questionTitle{51}{Chamada recursiva com acumulador}
\begin{flushright}
  \footnotesize
  \textit{\textbf{10: Ações de E/S recursivas}}\\
  \textit{c10-acumuladores-cauda-02}\\

\end{flushright}
Considere a versão com acumulador:

\begin{haskellcode}
lerESomar :: Integer -> IO Integer
lerESomar total = do
  n <- readLn
  if n == 0
    then return total
    else lerESomar (total + n)
\end{haskellcode}

Qual afirmação é a mais precisa?

\begin{enumerate}
  \item A função soma todos os valores apenas no caso base, sem atualizar o acumulador durante a recursão.

  \item O uso de acumulador garante, por si só, uso constante de memória em qualquer implementação Haskell.

  \item A chamada recursiva deixa de ser recursiva porque o acumulador substitui a função original.

  \item No ramo recursivo, a chamada \haskellinline|lerESomar (total + n)| é a última ação executada naquele ramo.

  \item A função passa a ser pura, pois o acumulador elimina o tipo \haskellinline|IO|.


\end{enumerate}

\questionTitle{52}{Associatividade da aplicação de função}
\begin{flushright}
  \footnotesize
  \textit{\textbf{3: Expressões e definições}}\\
  \textit{c03-aplicacao-notacoes-02}\\

\end{flushright}
A aplicação de função em Haskell associa à esquerda. Assim, como a expressão \haskellinline|logBase 2 8| deve ser lida?

\begin{enumerate}
  \item Como \haskellinline|(logBase 2 8)|, mas com \haskellinline|2| e \haskellinline|8| formando implicitamente uma tupla.

  \item Como \haskellinline|logBase (2, 8)|, pois funções de dois argumentos recebem tuplas por padrão.

  \item Como uma expressão inválida, pois \haskellinline|logBase| exige que os argumentos sejam separados por vírgula.

  \item Como \haskellinline|logBase (2 8)|, pois o argumento mais à direita é agrupado primeiro.

  \item Como \haskellinline|(logBase 2) 8|: primeiro obtém-se a função logarítmica de base \haskellinline|2|, depois ela é aplicada a \haskellinline|8|.


\end{enumerate}

\questionTitle{53}{Equivalência com guardas}
\begin{flushright}
  \footnotesize
  \textit{\textbf{5: Expressão condicional}}\\
  \textit{c05-if-aninhamento-03}\\

\end{flushright}
Qual definição com guardas é funcionalmente equivalente à função abaixo?
\begin{haskellcode}
sinal n = if n < 0 then -1 else if n > 0 then 1 else 0
\end{haskellcode}

\begin{enumerate}
  \item \begin{haskellcode}
sinal n
  | n > 0     = -1
  | n < 0     = 1
  | otherwise = 0
\end{haskellcode}

  \item \begin{haskellcode}
sinal n = if n == 0 then -1 else 1
\end{haskellcode}

  \item \begin{haskellcode}
sinal n
  | n < 0     = -1
  | otherwise = 1
\end{haskellcode}

  \item \begin{haskellcode}
sinal n
  | n <= 0    = -1
  | otherwise = 1
\end{haskellcode}

  \item \begin{haskellcode}
sinal n
  | n < 0     = -1
  | n > 0     = 1
  | otherwise = 0
\end{haskellcode}


\end{enumerate}

\questionTitle{54}{Classificação por faixas}
\begin{flushright}
  \footnotesize
  \textit{\textbf{5: Expressão condicional}}\\
  \textit{c05-guardas-avaliacao-03}\\

\end{flushright}
Analise a definição:
\begin{haskellcode}
conceito n
  | n >= 9    = "A"
  | n >= 7    = "B"
  | n >= 5    = "C"
  | otherwise = "D"
\end{haskellcode}
Qual é o resultado de \haskellinline|conceito 7.5|?

\begin{enumerate}
  \item \haskellinline|"D"|

  \item \haskellinline|"C"|

  \item \haskellinline|"A"|

  \item Ocorre um erro de tipos, pois as guardas usam números decimais.

  \item \haskellinline|"B"|


\end{enumerate}

\questionTitle{55}{Tipo de resultado de lookup}
\begin{flushright}
  \footnotesize
  \textit{\textbf{6: Estruturas de dados básicas}}\\
  \textit{c06-lookup-associacao-03}\\

\end{flushright}
Ao buscar uma chave em uma lista de associação com \haskellinline{lookup}, por que o resultado é representado com \haskellinline{Maybe}?

\begin{enumerate}
  \item Porque a chave pode estar presente ou ausente, e \haskellinline{Maybe} representa explicitamente essas duas possibilidades.

  \item Porque Haskell não permite devolver diretamente valores como \haskellinline{String} ou \haskellinline{Double}.

  \item Porque toda função que opera sobre tuplas em Haskell precisa devolver um \haskellinline{Maybe}.

  \item Porque o uso de \haskellinline{lookup} transforma automaticamente qualquer lista em uma enumeração.

  \item Porque \haskellinline{lookup} sempre devolve uma lista, mesmo quando encontra apenas um valor.


\end{enumerate}

\questionTitle{56}{Benefício de testar funções puras separadamente}
\begin{flushright}
  \footnotesize
  \textit{\textbf{10: Ações de E/S recursivas}}\\
  \textit{c10-separacao-puro-02}\\

\end{flushright}
Ao projetar um programa interativo em Haskell, por que é desejável separar a coleta de dados em \haskellinline|IO| do processamento puro?

\begin{enumerate}
  \item Porque funções puras imprimem seus resultados automaticamente no GHCi e no programa compilado.

  \item Porque o compilador exige que todo bloco \haskellinline|do| contenha exatamente uma função pura.

  \item Porque funções puras podem ser testadas e reutilizadas sem depender de entrada e saída reais.

  \item Porque ações de E/S não podem ter assinaturas de tipo explícitas.

  \item Porque funções puras são as únicas que podem chamar \haskellinline|readLn| com segurança.


\end{enumerate}

\questionTitle{57}{Sintaxe explícita e regra de layout}
\begin{flushright}
  \footnotesize
  \textit{\textbf{3: Expressões e definições}}\\
  \textit{c03-layout-avaliacao-02}\\

\end{flushright}
Qual trecho é equivalente à expressão com delimitadores explícitos \haskellinline|let { a = 1; b = 2 } in a + b|?

\begin{enumerate}
  \item \begin{haskellcode}
let a = 1
  b = 2
in a + b
\end{haskellcode}

  \item \begin{haskellcode}
let a = 1
in b = 2
in a + b
\end{haskellcode}

  \item \begin{haskellcode}
let a = 1
    b = 2
in a + b
\end{haskellcode}

  \item \begin{haskellcode}
a = 1
b = 2
in a + b
\end{haskellcode}

  \item \begin{haskellcode}
let a = 1; b = 2; a + b
\end{haskellcode}


\end{enumerate}

\questionTitle{58}{A saída segura do ambiente interativo}
\begin{flushright}
  \footnotesize
  \textit{\textbf{2: Ambiente de desenvolvimento}}\\
  \textit{c02-fluxo-trabalho-02}\\

\end{flushright}
Após terminar o desenvolvimento, testar exaustivamente o arquivo \texttt{.hs} e salvar o progresso, qual comando do GHCi é o responsável por encerrar a sessão interativa de forma limpa, devolvendo o controle ao terminal do sistema operacional?

\begin{enumerate}
  \item O comando \texttt{:close} (ou a abreviação \texttt{:c}).

  \item O comando \texttt{:quit} (ou a abreviação \texttt{:q}).

  \item O comando \texttt{:stop} (ou a abreviação \texttt{:s}).

  \item O comando \texttt{:exit} (ou a abreviação \texttt{:e}).

  \item O comando \texttt{:end} (sem abreviação suportada).


\end{enumerate}

\questionTitle{59}{Construção da lista durante o retorno da recursão}
\begin{flushright}
  \footnotesize
  \textit{\textbf{10: Ações de E/S recursivas}}\\
  \textit{c10-listas-ordem-02}\\

\end{flushright}
Analise a versão sem acumulador:

\begin{haskellcode}
lerLista :: IO [Integer]
lerLista = do
  x <- readLn
  if x == 0
    then return []
    else do
      resto <- lerLista
      return (x:resto)
\end{haskellcode}

Se o usuário digita \haskellinline|4|, depois \haskellinline|9|, e depois \haskellinline|0|, qual lista é produzida?

\begin{enumerate}
  \item \haskellinline|[0,9,4]|, pois a recursão reconstrói a lista de trás para frente incluindo o caso base.

  \item \haskellinline|[9,4]|, pois cada novo elemento é colocado no início da lista final.

  \item \haskellinline|[4,9]|, pois cada chamada adiciona seu valor à frente da lista produzida pelo restante da recursão.

  \item \haskellinline|[]|, pois \haskellinline|return []| descarta os valores anteriores.

  \item \haskellinline|[4,9,0]|, pois o sentinela também faz parte da lista retornada.


\end{enumerate}

\questionTitle{60}{Anatomia de uma definição em arquivo}
\begin{flushright}
  \footnotesize
  \textit{\textbf{3: Expressões e definições}}\\
  \textit{c03-arquivos-identificadores-02}\\

\end{flushright}
Considere a definição em um arquivo fonte:

\begin{haskellcode}
dobro n = 2 * n
\end{haskellcode}

Qual alternativa identifica corretamente seus componentes?

\begin{enumerate}
  \item A definição é inválida porque funções de um argumento exigem vírgula depois do nome.

  \item \haskellinline|n| é uma definição global independente criada automaticamente.

  \item \haskellinline|2 * n| é uma assinatura de tipo.

  \item \haskellinline|dobro| é o identificador definido, \haskellinline|n| é o parâmetro e \haskellinline|2 * n| é o corpo da definição.

  \item \haskellinline|dobro n| é o nome completo da função, e por isso ela não recebe argumentos.


\end{enumerate}

\questionTitle{61}{Equivalência entre if e guardas}
\begin{flushright}
  \footnotesize
  \textit{\textbf{5: Expressão condicional}}\\
  \textit{c05-if-vs-guardas-03}\\

\end{flushright}
Qual definição com guardas é equivalente à função abaixo?
\begin{haskellcode}
max2 x y = if x > y then x else y
\end{haskellcode}

\begin{enumerate}
  \item \begin{haskellcode}
max2 x y
  | x >= y    = y
  | otherwise = x
\end{haskellcode}

  \item \begin{haskellcode}
max2 x y
  | x < y     = x
  | otherwise = y
\end{haskellcode}

  \item \begin{haskellcode}
max2 x y
  | x > y     = True
  | otherwise = False
\end{haskellcode}

  \item \begin{haskellcode}
max2 x y
  | x > y     = x
  | otherwise = y
\end{haskellcode}

  \item \begin{haskellcode}
max2 x y = if x > y then y else x
\end{haskellcode}


\end{enumerate}

\questionTitle{62}{Assinatura em expressões}
\begin{flushright}
  \footnotesize
  \textit{\textbf{4: Tipos de Dados}}\\
  \textit{c04-assinatura-expressao-01}\\

\end{flushright}
Por que a expressão \haskellinline{read "42"} sozinha causaria um erro de ambiguidade de tipo, enquanto \haskellinline{read "42" :: Int} é aceita?

\begin{enumerate}
  \item A anotação é necessária porque \haskellinline{read} é uma função parcial.

  \item \haskellinline{read} é sobrecarregada e pode retornar vários tipos; a anotação \haskellinline{:: Int} resolve a ambiguidade.

  \item A anotação \haskellinline{:: Int} força a conversão da string para inteiro em tempo de execução.

  \item A anotação \haskellinline{:: Int} instrui o compilador a ignorar a inferência de tipos.

  \item \haskellinline{read} só funciona com anotação de tipo; sem ela, ocorre um erro de sintaxe.


\end{enumerate}

\questionTitle{63}{Otherwise fora da última guarda}
\begin{flushright}
  \footnotesize
  \textit{\textbf{5: Expressão condicional}}\\
  \textit{c05-otherwise-03}\\

\end{flushright}
Considere a definição:
\begin{haskellcode}
g x
  | otherwise = 1
  | x > 0     = 2
\end{haskellcode}
O que acontece ao avaliar \haskellinline|g 5|?

\begin{enumerate}
  \item O compilador reordena as guardas e produz \haskellinline|2|.

  \item O resultado é \haskellinline|1|, porque \haskellinline|otherwise| equivale a \haskellinline|True|.

  \item Ocorre um erro de sintaxe, pois \haskellinline|otherwise| só pode aparecer na última linha.

  \item O resultado é \haskellinline|2|, porque a guarda \haskellinline|x > 0| seria mais específica.

  \item Ocorre um erro em tempo de execução por ambiguidade entre as guardas.


\end{enumerate}

\questionTitle{64}{Tipo unitário}
\begin{flushright}
  \footnotesize
  \textit{\textbf{6: Estruturas de dados básicas}}\\
  \textit{c06-tuplas-unitario-01}\\

\end{flushright}
Qual afirmação descreve corretamente o tipo \haskellinline{()} em Haskell?

\begin{enumerate}
  \item É apenas outra notação para a lista vazia \haskellinline{[]}.

  \item É um tipo reservado exclusivamente para expressões numéricas.

  \item É um sinônimo de tipo para \haskellinline{Bool}.

  \item É a forma como Haskell representa uma tupla de um elemento.

  \item É o tipo unitário: possui exatamente um único valor, também escrito como \haskellinline{()}.


\end{enumerate}

\questionTitle{65}{Escolha entre Int e Integer}
\begin{flushright}
  \footnotesize
  \textit{\textbf{9: Funções recursivas}}\\
  \textit{c09-eficiencia-integer-int-01}\\

\end{flushright}
Por que o capítulo usa \haskellinline|Integer|, e não \haskellinline|Int|, em exemplos como \haskellinline|fatorial|?

\begin{enumerate}
  \item Porque \haskellinline|Int| não permite valores positivos.

  \item Porque \haskellinline|Integer| tem precisão arbitrária, enquanto \haskellinline|Int| tem tamanho fixo e pode sofrer estouro de capacidade.

  \item Porque \haskellinline|Integer| torna qualquer função automaticamente recursiva de cauda.

  \item Porque \haskellinline|Integer| impede chamadas recursivas infinitas.

  \item Porque \haskellinline|Int| só pode ser usado em listas.


\end{enumerate}

\questionTitle{66}{Compatibilidade dos ramos}
\begin{flushright}
  \footnotesize
  \textit{\textbf{5: Expressão condicional}}\\
  \textit{c05-if-regras-03}\\

\end{flushright}
Qual expressão abaixo é rejeitada porque os ramos \haskellinline|then| e \haskellinline|else| não produzem valores do mesmo tipo?

\begin{enumerate}
  \item \haskellinline|if odd 5 then "ímpar" else "par"|

  \item \haskellinline|if odd 5 then 1 else 0|

  \item \haskellinline|if odd 5 then 'i' else 'p'|

  \item \haskellinline|if odd 5 then "ímpar" else 0|

  \item \haskellinline|if odd 5 then [1,2] else [3,4]|


\end{enumerate}

\questionTitle{67}{Recursão múltipla em Fibonacci}
\begin{flushright}
  \footnotesize
  \textit{\textbf{9: Funções recursivas}}\\
  \textit{c09-formas-fibonacci-02}\\

\end{flushright}
Na definição
\begin{haskellcode}
fib n
  | n == 0 = 0
  | n == 1 = 1
  | n > 1  = fib (n - 2) + fib (n - 1)
\end{haskellcode}
por que dizemos que há recursão múltipla?

\begin{enumerate}
  \item Porque o caso recursivo realiza mais de uma chamada à própria função.

  \item Porque a função retorna uma lista com vários elementos.

  \item Porque a função usa duas guardas antes do caso recursivo.

  \item Porque a função chama outra função auxiliar com acumuladores.

  \item Porque há mais de um caso base.


\end{enumerate}

\questionTitle{68}{Avaliação com where e dependência local}
\begin{flushright}
  \footnotesize
  \textit{\textbf{3: Expressões e definições}}\\
  \textit{c03-definicoes-where-02}\\

\end{flushright}
Considere a definição:

\begin{haskellcode}
f x = a + b
  where
    a = x + 1
    b = a * 2
\end{haskellcode}

Qual é o valor de \haskellinline|f 3|?

\begin{enumerate}
  \item Ocorre erro, pois \haskellinline|b| não pode usar \haskellinline|a| definido no mesmo bloco \haskellinline|where|.

  \item \haskellinline|8|.

  \item \haskellinline|12|.

  \item \haskellinline|7|.

  \item \haskellinline|4|.


\end{enumerate}

\questionTitle{69}{sort como combinação de estrutura paramétrica e restrição}
\begin{flushright}
  \footnotesize
  \textit{\textbf{7: Polimorfismo}}\\
  \textit{cap07-comparacao-sort-param-adhoc}\\

\end{flushright}
A assinatura abaixo combina duas ideias:

\begin{minted}{haskell}
sort :: Ord a => [a] -> [a]
\end{minted}

Qual interpretação é mais adequada?

\begin{enumerate}
  \item A função é uma conversão numérica, pois \haskellinline{Ord} pertence apenas à hierarquia de \haskellinline{Num}.
  \item A presença de \haskellinline{Ord a} torna a função monomórfica, aceitando apenas \haskellinline{Int}.
  \item A função depende de herança entre objetos mutáveis.
  \item A estrutura da função é polimórfica sobre listas de \haskellinline{a}, mas a ordenação dos elementos exige uma instância \haskellinline{Ord a}.
  \item A assinatura mostra polimorfismo paramétrico puro, pois classes de tipo não impõem restrições reais.

\end{enumerate}


\end{multicols}
\makeanswersheet{UFOP}{BCC222: Programação Funcional}{José Romildo}{2026/1}{Primeira Prova}{06}{69}{25}{7}
\gabarito{06}{E,E,E,C,A,E,C,C,C,D,C,B,E,E,B,A,D,C,E,B,C,A,C,E,B,E,B,B,B,A,E,B,D,B,B,D,D,D,A,D,D,C,E,C,A,A,B,B,D,D,D,E,E,E,A,C,C,B,C,D,D,B,B,E,B,D,A,C,D}
\end{document}
